% Compile twice: lualatex 20001318.tex
\documentclass[11pt]{article}
\usepackage[margin=25mm]{geometry}
\usepackage{fontspec,unicode-math}
\setmainfont{Latin Modern Roman}
\setmathfont{Latin Modern Math}
\usepackage{amsmath,amssymb,mathtools}
\usepackage{xurl,hyperref}
\hypersetup{colorlinks=true,urlcolor=blue,linkcolor=blue}
\setcounter{secnumdepth}{0}
\setlength{\parindent}{0pt}
\setlength{\parskip}{5pt}
\setlength{\emergencystretch}{3em}
\newcommand{\sref}[2]{\hyperlink{source:#1:#2}{[S#2]}}
\newcommand{\cataloguescope}[1]{\paragraph{Source question.}#1}
\newcommand{\floor}[1]{\left\lfloor#1\right\rfloor}
\begin{document}
% UnsolvedMath ID 20001318; source catalogue code AIM-CRYPTOGRAPHY-0005
\setcounter{section}{20001317}
\section{The inverse-state twist interface, extraction boundary, and exact storage limit}\label{20001318}
{\small\sffamily UnsolvedMath ID 20001318 · Computer Science}
\subsection{Definitions and mathematical statement}

\paragraph{Shared notation.}Standard number systems are denoted by $\mathbb N,\mathbb Z,\mathbb Q,\mathbb R,\mathbb C$; index ranges, computational models and specialized notation follow the cited problem source and the definitions below.


A digital signature enables public verification of an authenticated message. Ring signatures hide which member of a stated group signed; blind signatures conceal the signed message from the signer; threshold signatures distribute signing authority. The source asks for constructions from group actions and improved memory use through an analogue of the twist technique.
\paragraph{Statement recorded in the supplied source.}
Advanced digital signature functionality from group theory

The connection between SDLP and group actions suggests the availability of generic group-action tools to build more efficient signatures, or signatures with more specialised properties.

We would like to use generic group-action tools to build ring signatures, blind signatures, threshold signatures and more. We would also like to apply the "twist" technique present in similar isogeny schemes to improve the memory performance of group-based digital signatures. \sref{20001318}{2}
\paragraph{Scope qualification.}
The source names functionalities but supplies no security experiment, group family, or formal description of the proposed twist. This is a construction-method question, not a claimed secure protocol.
\paragraph{Source status.}
The supplied catalogue labels this record partially solved. Its source-status metadata was checked against the supplied snapshot on 2026-10-11. This import does not independently establish current open status or mathematical correctness.
\subsection{Short English statement}
Can group-action methods build specialized digital signatures and reduce their memory requirements?
\subsection{Sources}
\begin{itemize}
\item[\hypertarget{source:20001318:1}{S1}] Ulam AI and the UnsolvedMath Contributors, UnsolvedMath: A Curated Collection of Open Mathematics Problems, record 20001318 (AIM-CRYPTOGRAPHY-0005). Catalogue curation and metadata: CC BY 4.0. The supplied source question is retained; definitions and the English summary are editorial additions.
\url{https://huggingface.co/datasets/ulamai/UnsolvedMath}.
\item[\hypertarget{source:20001318:2}{S2}] AIM Workshop Problem Lists, original problem source linked by catalogue record AIM-CRYPTOGRAPHY-0005
\url{http://aimpl.org/postquantgroup/2/}.
\end{itemize}
\paragraph{Preparation and review.}
Original question attributed to the cited source and UnsolvedMath contributors. Codex (GPT-6; exact serving revision not exposed) assisted with definitions, English summaries, source inspection and formatting. No new human mathematical authorship or independent proof verification is asserted. Canonical ID 20001318 and source code AIM-CRYPTOGRAPHY-0005 checked against the complete supplied registry and existing repository statements; known flagged repetitions and obvious identical conjectures excluded. Full mathematical-equivalence/literature review has not been performed. Statement preparation and final review followed contribution prompts C and D; automated checks verify structure and identity, not open status.
\label{end:20001318}
\end{document}
